
\documentclass[10pt,xetex]{beamer} %[compress, blue]
\mode<presentation>


\setbeamercolor{frametitle}{fg=lightred,bg=darkred}
\setbeamercolor{title}{fg=lightred,bg=darkred}

\usepackage[russian,english]{babel}

\usepackage{multirow}
\usepackage{fontspec}
\usepackage{xunicode}
\usepackage{xltxtra}
\usepackage{wasysym}
\usepackage{url}
\defaultfontfeatures{Mapping=tex-text} %PunctuationSpace=3,Scale=MatchLowercase,,Numbers=Lining}
\setmainfont[]{Nimbus Sans L}
\setsansfont[]{Nimbus Sans L}
\setmonofont[]{Liberation Mono}

% Chuvash colours are 203,0,0 and 255,254,0
\definecolor{darkred}{RGB}{153,0,0}
\definecolor{lightred}{RGB}{226,200,200}
\definecolor{yellow}{RGB}{226,226,50}


\usetheme{Warsaw}

\usecolortheme[named=darkred]{structure}


\newfontfamily\cyrtext[Scale=MatchLowercase]{Nimbus Sans L}
\useoutertheme[subsection=false]{smoothbars}
%\setbeamertemplate{footline}[page number]{} % header empty
% \setbeamercolor{footer}{fg=blue,bg=lightred}
%\setbeamercolor{footline}{fg=white,bg=lightred} 
%    \setbeamercolor{page number in head/foot}{use=headline,fg=headline.fg,bg=frametitle.fg}
\setbeamertemplate{footline}[page number]{} % header empty

\setbeamertemplate{headline}[text line]{} % header empty

\setbeamertemplate{navigation symbols}{}

\newcommand{\atag}[1]{{\scriptsize{\texttt{#1}}}}
\newcommand{\itemise}{\itemize}

%\date{20th May 2012}
\date{}
\title{A prototype machine translation system for Tatar and Bashkir \\
              based on free/open-source components}

%\author{Francis M. Tyers \and Jonathan North Washington \\ 
%       U.A \and I.U \\
%Ilnar Salimzyanov \and Rustam Battalov}
%
%\author{{\bf Francis M. Tyers}\\ {\tt \href{mailto:ftyers@dlsi.ua.es}{ftyers@dlsi.ua.es}}\\Dept. de Lleng. i Sist. Inform.,\\Universitat d'Alacant \and
%{\bf Jonathan North Washington}\\ {\tt \href{mailto:jonwashi@indiana.edu}{jonwashi@indiana.edu}}\\Department of Linguistics,\\Indiana University\\
%{\bf Ilnar Salimzyanov} \and
%{\bf Rustam Battalov} 
%}
\begin{document}

\begin{frame} % FIXME: Get rid of the space between the title and the names etc.
        \titlepage
\centering
\begin{tabular}{cc}
 {\bf Francis M. Tyers} & {\bf Jonathan North Washington} \\
 Dept. Lleng. i Sist. Inform. & Dept. Linguistics \\
 Universitat d'Alacant & Indiana University \\
 {\tt \href{mailto:ftyers@dlsi.ua.es}{ftyers@dlsi.ua.es}} & {\tt \href{mailto:jonwashi@indiana.edu}{jonwashi@indiana.edu}} \\ 
 ~ & ~ \\
 {\bf Ilnar Salimzyanov} & {\bf Rustam Batalov} \\
 Kazan Federal University & Bashkir State University \\
 {\tt \href{mailto:ilnar.salimzyan@gmail.com}{ilnar.salimzyan@gmail.com}} & {\tt \href{mailto:taqmaq@mail.ru}{taqmaq@mail.ru}} \\
 ~ & ~ \\
 ~ & ~ \\
\end{tabular}

\end{frame}

%%%%%%%%%%%%%%%%%%%%%%%%%%%%%%%%%%%%%%%%%%%%%%%%%%%%%%%%%%%%%%%%%%%%%%%%%%%%%%%
\begin{frame}
  \frametitle{Introduction}
  
  \begin{onlyenv}<1>
    \begin{block}{Layout of the talk}
    
      \begin{itemize}
        \item Introduction
        \item Description of architecture and development process
        \item Evaluation
        \item Discussion and future work
      \end{itemize}
    
    \end{block}
  \end{onlyenv}
  
  \begin{onlyenv}<2>
    \begin{block}{Why make a Tatar $\leftrightarrow$ Bashkir MT system?}
    
      \begin{itemize}
        \item There are quite a few examples of MT systems made with Apertium
        \begin{itemize}
          \item but no working examples between ``morphologically complex'' languages
        \end{itemize}
    
        \item Course in Šupaškar (Čeboksary), January 2012
        \begin{itemize}
          \item Nearly all the languages of Russia are ``morphologically complex'', but
              we had no example system to present.
        \end{itemize} 
    
        \item The idea was to create a pedagogical example
        \begin{itemize}
          \item A simple demonstration system 
          \item that nonetheless showed all modules working
        \end{itemize}
     
        \item We chose Tatar and Bashkir because,
        \begin{itemize}
          \item Native speakers available to help out
          \item Good results could be accomplished with little work
        \end{itemize}
      \end{itemize}
    \end{block}
  \end{onlyenv}
  
\end{frame}

\begin{frame}
  \frametitle{Tatar and Bashkir} % Map
  
  \begin{onlyenv}<1>
    \includegraphics[width=\textwidth]{middlevolga19001500.jpg}
  \end{onlyenv}
  \begin{onlyenv}<2>
    
    \begin{block}{Tatar}
      
      \begin{itemize}
        \item spoken by $>$ 6.5m  people
        \item hargle bargle %FIXME: more real stuff needed
      \end{itemize}
      
    \end{block}
      
    \begin{block}{Bashkir}
      
      \begin{itemize}
        \item spoken by $>$ 1.3m people
        \item some facts we made up %FIXME: put real stuff here
      \end{itemize}
    
    \end{block}
    
  \end{onlyenv}
  \begin{onlyenv}<3>
  
    %FIXME: Add IPA 
    \begin{block}{Tatar and Bashkir}
      Һава бүген бик әйбәт, җылы гына.  \\
      Һауа бөгөн бик әйбәт, йылы ғына. \\
    \end{block}
    
    \begin{block}{Comparison with other Turkic languages}
      Hava bugün çok güzel, yeterince sıcak. \\
      Ҫанталӑк паян питӗ хитре, самай ӑшӑ. \\
      Ауа райы бүгін әбден жақсы, жылы. \\
      Аба бүгүн абдан жакшы, жетишинче жылуу. 
    \end{block}
  
  \end{onlyenv}
  
\end{frame}

\begin{frame}
  \frametitle{Apertium}

\begin{center}
  \includegraphics[width=0.82\textwidth]{apertiumlangs.pdf}%FIXME: svg->pdf is ugly
\end{center}
\begin{block}{Overview}
  
  \begin{itemize}
    \item Free/open-source software (GPL licence)
    \begin{itemize}
      \item Not only the programs, but the data too!
    \end{itemize}
    \item Fairly mature data for over 30 language pairs
    \item ... but so far no mature data for Turkic languages!
  \end{itemize}
  
\end{block}


\end{frame}

%%%%%%%%%%%%%%%%%%%%%%%%%%%%%%%%%%%%%%%%%%%%%%%%%%%%%%%%%%%%%%%%%%%%%%%%%%%%%%%
\begin{frame}
  \frametitle{Architecture}

\includegraphics[width=\textwidth]{../architecture.pdf}

\end{frame}

\begin{frame}
  \frametitle{Tools}

  \begin{block}{Helsinki Finite-State Toolkit (HFST)}
    \begin{itemize}
      \item Free/open-source implementation of the Xerox 
         finite-state formalisms {\small {\tt lexc}}/{\small {\tt twol}}
      \item Used for morphological analysis and generation
    \end{itemize}
  \end{block}
  
  \begin{block}{VISL Constraint Grammar}
    \begin{itemize}
      \item Rule-based morphological disambiguation
    \end{itemize}
  \end{block}
  
  \begin{block}{Apertium}
    \begin{itemize}
      \item Lexical and structural transfer
      \item ``Platform'': Build system, ancilliary tools, etc.
    \end{itemize}
  \end{block}
  
\end{frame}


\begin{frame}[fragile]

\fontsize{9pt}{11.2}\selectfont

\begin{tabular}{l}
%\hline
%{\bf Stage} & {\bf Representation} \\
\hline
\hline
 Һава бүген бик әйбәт, җылы гына. \\ 
\hline
{\bf Morphological analysis} \\
\hline
 \^{}Һава/һава\atag{<n>}\atag{<attr>}/һава\atag{<n>}\atag{<nom>\$} \^{}бүген/бүген\atag{<adv>\$}  \\
              \^{}бик/бик\atag{<adv>}/бик\atag{<n>}\atag{<attr>}/бик\atag{<n>}\atag{<nom>\$} \\ 
              \^{}әйбәт/әйбәт\atag{<adj>}/әйбәт\atag{<adj>}\atag{<subst>}\atag{<nom>\$}\^{},/,\atag{<cm>\$} \\ 
              \^{}җылы/җылы\atag{<n>}\atag{<attr>}/җылы\atag{<n>}\atag{<nom>}/җылы\atag{<adj>}/җылы\atag{<adj>}\atag{<subst>}\atag{<nom>\$} \\
              \^{}гына/гына\atag{<postadv>\$}\^{}./.\atag{<sent>\$} \\
\hline
{\bf Morphological disambiguation} \\
\hline
 \^{}Һава\atag{<n>}\atag{<nom>\$} \^{}бүген\atag{<adv>\$} \^{}бик\atag{<adv>\$} \^{}әйбәт\atag{<adj>\$}\^{},\atag{<cm>\$} \\ 
              \^{}җылы\atag{<adj>\$} \^{}гына\atag{<postadv>\$}\^{}.\atag{<sent>\$}\\
\hline
{\bf Lexical transfer and selection} \\
\hline
 \^{}Һава\atag{<n>}\atag{<nom>}/Һауа\atag{<n>}\atag{<nom>\$} \^{}бүген\atag{<adv>}/бөгөн\atag{<adv>\$} \^{}бик\atag{<adv>}/бик\atag{<adv>\$} \\
              \^{}әйбәт\atag{<adj>}/әйбәт\atag{<adj>\$}\^{},\atag{<cm>}/,\atag{<cm>\$} \^{}җылы\atag{<adj>}/йылы\atag{<adj>\$} \\ 
              \^{}гына\atag{<postadv>}/ғына\atag{<postadv>\$}\^{}.\atag{<sent>}/.\atag{<sent>\$}\\
\hline
{\bf Structural transfer} \\
\hline
 \^{}Һауа\atag{<n>}\atag{<nom>\$} \^{}бөгөн\atag{<adv>\$} \^{}бик\atag{<adv>\$} \^{}әйбәт\atag{<adj>\$}\^{},\atag{<cm>\$} \\ 
              \^{}йылы\atag{<adj>\$} \^{}ғына\atag{<postadv>\$}\^{}.\atag{<sent>\$}\\
\hline
{\bf Morphological generation} \\
\hline
 Һауа бөгөн бик әйбәт, йылы ғына. \\
\hline
\hline
\end{tabular}

\end{frame}

%%%%%%%%%%%%%%%%%%%%%%%%%%%%%%%%%%%%%%%%%%%%%%%%%%%%%%%%%%%%%%%%%%%%%%%%%%%%%%%
\begin{frame}
  \frametitle{Development process}

  \begin{itemize} % FIXME: check this
    \item Stems in both Tatar and Bashkir are added to an online spreadsheet
    \begin{itemize}
      \item Added according to frequency in the RNC. % downsides
      \item At the same time as translations, POS categories are added, and some
        relevant subcategorisations (transitivity, Tatar infinitive)
    \end{itemize}
    \item Scripts are used to automate the process of converting the spreadsheet
       to bilingual ({\small {\tt .dix}}) and monolingual ({\small {\tt .lexc}}) formats.
    \item At the same time, the morphotactics and phonology are written
       according to available grammatical descriptions
    
  \end{itemize}
\end{frame}

\begin{frame}
  \frametitle{Speed}

\begin{center}
\begin{tabular}{cc}

\includegraphics[angle=270,width=0.5\textwidth]{hist-tt.ps} &  \includegraphics[angle=270,width=0.5\textwidth]{hist-ba.ps} \\
Tatar & Bashkir \\
\end{tabular}
\end{center}

\end{frame}

%%%%%%%%%%%%%%%%%%%%%%%%%%%%%%%%%%%%%%%%%%%%%%%%%%%%%%%%%%%%%%%%%%%%%%%%%%%%%%%
\begin{frame}
  \frametitle{Evaluation}

\begin{block}{Rules and lexica}
  \begin{center}
    \begin{tabular}{lr}
    \hline
     {\bf Corpus}                        & {\bf Entries}\\
    \hline
     Bilingual dictionary                & 2,685     \\
     Disambiguation ({\small {\tt tt}} and {\small {\tt ba}}) & 6   \\
     Transfer ({\small {\tt tt$\rightarrow$ba}}) & 3 \\
     Transfer ({\small {\tt ba$\rightarrow$tt}}) & 3 \\
    \hline
    \end{tabular}
  \end{center}  

\end{block}

\begin{block}{Coverage}

\begin{center}
  \begin{tabular}{lrr}
  \hline
   {\bf Corpus}             & {\bf Tokens}    & {\bf Coverage}\\
  \hline
   Tatar New Testament (NT) & 163,603   & 72.04\% \\
   Tatar Wikipedia          & 37,123    & 70.19\% \\
  \hline
   Bashkir Wikipedia        & 12,267    & 65.99\% \\
  \hline
  \end{tabular}

\end{center}
\end{block}

\begin{block}{Error rate}
\begin{center}
  % NOTE: The performance without MT is:
  %        tt->ba: 56,59% 
  %        ba->tt: 56,41% 
  % simple 'transliteration' (e.g. ҙ ҡ ғ → з к г):
  %        ba->tt: 46.47%
  %
  \begin{tabular}{ccrrr} 
  \hline
   {\bf Corpus}           & {\bf Direction}   & {\bf Tokens}  & {\bf Unknown} & {\bf WER}  \\
  \hline
   \multirow{2}{*}{story} & {\small {\tt tt$\rightarrow$ba}} & 311     & 2  & 6.73\% \\
                          & {\small {\tt ba$\rightarrow$tt}} & 312     & 0  & 6.43\%  \\
  \hline
  \end{tabular}
\end{center}
\end{block}

\end{frame}

\begin{frame}
  \frametitle{Weak points}

\begin{block}{Linguistic}
  
  \begin{itemize} % FIXME: check if these are still true
    \item Majority of errors due to mistakes and gaps in Tatar morphophonology
    \item Combinations of case and possessive suffixes
    \item Orthographical representations of phonology
    \item Coverage
    \item Hargle bargle
  \end{itemize}
  
\end{block}

\begin{block}{Technical}
  \begin{itemize}
    \item Vowel harmony processing on clitics (e.g., \emph{да}/\emph{дә} `and')  
       after unknown words.
    \item Traditional methods of quality control are difficult to apply
  \end{itemize} 
\end{block}

\end{frame}

\begin{frame}
  \frametitle{Observations}

\begin{block}{Consistency consistency}
  \begin{itemize}
      \item For tagsets: If something is the same, tag it the same
      \item Otherwise you spend a lot of time on useless `transfer'
        (e.g. ky, ~50 transfer rules, around 4 `real' rules)
      \item If you solve a problem well in one language, and it comes up 
        in another one, solve it the same way (e.g. epenthesis, vowel harmony)
  \end{itemize}
\end{block}

\begin{block}{Parallel development}
  \begin{itemize}
    \item Building the transducers and bilingual dictionary at the same time is good
  \end{itemize}
\end{block}

\begin{block}{Starting from scratch}
  \begin{itemize}
    \item Much easier to work with new code than adapt existing code
    \begin{itemize}
      \item Not to say that existing code can't be useful as a model 
    \end{itemize}
  \end{itemize}
\end{block}

\end{frame}

%%%%%%%%%%%%%%%%%%%%%%%%%%%%%%%%%%%%%%%%%%%%%%%%%%%%%%%%%%%%%%%%%%%%%%%%%%%%%%%

\begin{frame}
  \frametitle{Future work}

\begin{block}{For Tatar$\leftrightarrow$Bashkir:}
  \begin{itemize}
    \item Expand lexicons 
    \item Improve morphophonology
  \end{itemize}
\end{block}

\begin{block}{Other Turkic language projects:}
  \begin{itemize}
    \item We're currently working on:

    \begin{itemize}
      \item Turkish$\leftrightarrow$Chuvash
      \item Tatar$\leftrightarrow$Kazakh$\dag$
      \item Turkish$\leftrightarrow$Turkmen$\dag$
      \item Turkish$\leftrightarrow$Tatar$\dag$
    \end{itemize}

    \item We have worked on: 
    \begin{itemize}
      \item Turkish$\leftrightarrow$Azerbaijani$\ddag$
      \item Turkish$\leftrightarrow$Kyrgyz$\ddag$
    \end{itemize}
  \end{itemize}
  $\dag$ As projects in the 2012 Google Summer of Code. \\
  $\ddag$ As projects in the 2011 Google Summer of Code.
\end{block}

\end{frame}

\begin{frame}
  \frametitle{Conclusion}

\begin{block} % FIXME: Write a better conclusion

  \begin{itemize}
    \item We've presented
    \begin{itemize}
      \item a prototype MT system for translating between Tatar and Bashkir,
      \item described the tools used to make it,
      \item given a preliminary evaluation of the translation quality,
      \item and given some observations about the development process
    \end{itemize}
  \end{itemize}


\end{block}

\begin{center}
Try it out! \url{http://elx.dlsi.ua.es/~fran/tt-ba/}
\end{center}

\end{frame}

\begin{frame}

\begin{Large}
\begin{center}
Рәхмәт \\ 
Räxmät \\
Рахмәт \\ 
Teşekkürler\\
~\\
\smiley
\end{center}
\end{Large}
\end{frame}


\end{document}
